\centerline{\bf \Huge Окр + $H$}
\centerline{\bf \Huge }


\problem{\textbf{1}} Четырёхугольник $ABCD$ вписан в окружность, $AC$ — ее диаметр. Пусть $K$ и $M$ — проекции вершин $A, C$ на диагональ $BD$ соответственно. Через точку $K$ проведена прямая, параллельная $BC$ и пересекающая $AC$ в точке $P$. Докажите, что $\angle KMP = 90^{\circ}◦$

\problem{\textbf{2}} Докажите, что окружность, проходящая через основания проекций вершины $D$ параллелограмма $ABCD$ на стороны треугольника $ABC$, проходит через точку пересечения диагоналей
параллелограмма.

\problem{\textbf{3}} В треугольнике $ABC$ на стороне $AC$ нашлись такие точки $D$ и $E$, что $AB = AD$ и
$BE = EC$. При этом точка $E$ лежит между точками $A$ и $D$. Точка $F$ — середина дуги $BC$
окружности, описанной около треугольника $ABC$. Докажите, что точки $B, E, D, F$ лежат на
одной окружности.

\problem{\textbf{4}} Докажите, что отличная от $B$ точка пересечения окружностей, построенных на сторонах
$BA$ и $BC$ треугольника $ABC$ как на диаметрах, лежит на прямой $AC$.

\problem{\textbf{5}} В треугольнике $ABC$ сторона $AC$ наименьшая. На сторонах $AB$ и $BC$ взяты точки $K$
и $L$ соответственно, причём $KA = AC = CL$. Пусть $M$ — точка пересечения $AL$ и $KC$, а $I$ —
точка пересечения биссектрис треугольника $ABC$. Докажите, что прямая $MI$ перпендикулярна
прямой $AC$.

\problem{\textbf{6}} В остроугольном треугольнике $ABC$ провели высоту $CC_1$ и медиану $AA_1$. Оказалось,
что точки $A$, $A_1$, $C$, $C_1$ лежат на одной окружности.

а) Докажите, что треугольник $ABC$ равнобедренный.

б) Найдите площадь треугольника $ABC$, если $AA_1 : CC_1$ = 5 : 4 и $A_1C_1$ = 4.

\problem{\textbf{7}} В остроугольном треугольнике $ABC$ $\angle A = 60^{\circ}◦$
. Высоты $BN$ и $CM$ треугольника $ABC$
пересекаются в точке $H$. Точка $O$ — центр окружности, описанной около треугольника $ABC$.

а) Докажите, что $AH$ = $AO$.

б) Найдите площадь треугольника $AHO$, если $BC$ = $6\sqrt{3}$, $\angle ABC = 45^{\circ}◦$
